\section{Oversampling on the original fixed window}
\label{sec:compact}
The essential change is to interpolate to an order $N$ \emph{larger} than the
jet order $R$. Shrinking the physical grid to suppress an order-$R$ remainder
unnecessarily pays the inverse condition number again at every derivative.
Instead we keep the window fixed and suppress the remainder by oversampling.

\begin{lemma}[Exponential interpolation bound]
\label{lem:grid}
Let $V_N=(k^r/r!)_{0\le k,r\le N}$. Then
$\|V_N^{-1}\|_\infty\le2^N\binom{2N}{N}\le8^N$.
For $f=h_\beta-h_\gamma$, $N\ge R\ge1$ and $0<u\le1$, put
\[
 A_{N,R,u}=8^N u^{-R},\qquad
 E_{N,u}=2e^{\Lambda Nu}\frac{(\Lambda Nu)^{N+1}}{(N+1)!}.
\]
Then
\begin{equation}
 \max_{r\le R}|f^{(r)}(0)|\le
 A_{N,R,u}\left(\max_{1\le k\le N}|f(ku)|+E_{N,u}\right).
 \label{eq:oversample}
\end{equation}
The zero node is supplied by $f(0)=0$ and requires no observation.
\end{lemma}
\begin{proof}
The $k$th Lagrange polynomial is
$L_k(x)=\prod_{l\ne k}(x-l)/\prod_{l\ne k}(k-l)$.
An inverse-matrix entry is $r![x^r]L_k(x)$. Since $r!\le N^r$, its absolute
value is at most the absolute-coefficient numerator evaluated at $N$, divided
by $k!(N-k)!$. This is at most
$((2N)!/N!)/(k!(N-k)!)$. Summing over $k$ gives
$2^N\binom{2N}{N}$. The usual central-binomial bound gives $8^N$.
Taylor's formula and $|f^{(N+1)}(t)|\le2e^{\Lambda t}\Lambda^{N+1}$ give a
remainder at most $E_{N,u}$ at all grid nodes. Invert $V_N$ for the scaled
jets $u^r f^{(r)}(0)$ and use $u^{-r}\le u^{-R}$ for $r\le R$.
\end{proof}

For $j=J+1$, $R=4J+8$, $L=Q^{25j}$ and $v=\log(1/\delta)$,
$0<\delta\le1$, set $\tau=T/4$ and define the multiple of four
\begin{equation}
\begin{split}
 H_c&=\max\left\{8,2R,T/2,256e^2\Lambda^2\tau,
 \frac{\log4+\Lambda\tau+25j\log Q+v}{\log2}\right\},\\
 N_c&=4\lceil H_c/4\rceil,\qquad
 A_c=8^{N_c}(N_c/\tau)^R,\qquad
 p_c=2^{-N_c/4-1}/N_c,\\
 \kappa_c(J,\delta)&=\frac{p_c\delta^2}{4L^2 A_c^2}.
\end{split}\label{eq:compact-certificate}
\end{equation}
This uses the existing cells $(K,1,k)$ with $K=(N_c-8)/4$: because
$N_c\ge T/2$, their times are precisely $k\tau/N_c$ and their masses are
$p_c$. Thus it does not introduce a new sampling distribution.

\begin{theorem}[Fixed-window separation with linear-logarithmic depth cost]
\label{thm:compact}
Define the explicit constants
\begin{align}
 K_c&=32+T+256e^2\Lambda^2\tau+2\Lambda\tau+50\log Q,\notag\\
 C_c&=32+50\log Q+5K_c+17\log K_c+16|\log\tau|.
 \label{eq:compact-constants}
\end{align}
For all $J\ge0$, $0<\delta\le1$, and $d_J(\beta,\gamma)\ge\delta$,
\begin{equation}
 D_\omega(\beta,\gamma)\ge\kappa_c(J,\delta)
 \ge e^{-C_cj\log(ej)}\delta^{12},\qquad j=J+1.
 \label{eq:compact-separation}
\end{equation}
\end{theorem}
\begin{proof}
For $N=N_c$, $u=\tau/N$, use $N\ge2R$, $N/\tau\ge1$, and
$(N+1)!\ge((N+1)/e)^{N+1}$. Then
\[
 A_c E_{N,u}\le
 2e^{\Lambda\tau}
 \left(8e\Lambda\sqrt{\tau/N}\right)^N\frac{e\Lambda\tau}{N}
 \le 2e^{\Lambda\tau}2^{-N}.
\]
The definition of $N$ ensures $LA_cE_{N,u}\le\delta/2$.
Equations~\eqref{eq:jet} and \eqref{eq:oversample} force some positive grid
value to have magnitude at least $\delta/(2LA_c)$. Its weighted square proves
the first assertion.

We give the complete logarithmic reduction. From the maximum defining $H_c$,
rounding by at most four, $R\le8j$, and $1/\log2<2$, one obtains
$N\le K_cj+2v$. Put $a_0=2\log8+(\log2)/4<5$. Directly from the certificate,
\[
 -\log\kappa_c=2v+\log8+50j\log Q+a_0N+2R\log(N/\tau)+\log N.
\]
Since $\log N\le\log(K_cj)+2v/(K_cj)$ and $K_c\ge32$, the coefficient
of $v$ on the right is at most
$2+2a_0+34/K_c<12$. The remaining terms are at most
\[
 j\{\log8+50\log Q+5K_c+17\log K_c+16|\log\tau|\}
       +17j\log j\le C_cj\log(ej).
\]
This includes $j=1$, arbitrary positive $T$, and arbitrary $\delta\le1$.
\end{proof}

\begin{remark}
The report's $e^{-1000Fj^2}\delta^{20j}$ certificate already improves the
preceding cubic-depth bound. Theorem~\ref{thm:compact} strengthens its depth
\emph{order} by separating interpolation order from reconstruction depth.
No inference is made that a previously valid weaker sufficient theorem was
false. The constants here are deliberately explicit, not numerically optimized.
\end{remark}

